\section{Properties of the Coordination Interface}
\label{sec:properties}
The interface admits several limited statements that help interpret the experiments. They concern the meaning of the executed commands and comparisons, rather than optimality of the learned policy.

\subsection{Cumulative Feasibility of a Trip Budget}
Suppose a trip first receives allowance $B_0\geq0$. Let $h_1,\ldots,h_m$ be its subsequently committed holding commands and let $u_1,\ldots,u_q\geq0$ be the reductions imposed by later budget publications. If each command is constrained by the currently remaining allowance and debited exactly once, then
\begin{equation}
 B_{\mathrm{rem}}=B_0-\sum_{e=1}^{m}h_e-\sum_{v=1}^{q}u_v\geq0.
 \label{eq:budget_identity}
\end{equation}
This follows by induction over the ordered sequence of command submissions and publications. A legal submission subtracts at most the current remainder. A publication replaces the remainder by the minimum of itself and a nonnegative proposal, subtracting a nonnegative reduction. Neither event increases the allowance. Consequently, total committed holding after issuance cannot exceed $B_0$, regardless of how often the upper mode changes.

The same identity holds inside an individual predictive continuation because its ledger uses the same events. It does not imply a realized return-time bound: uncertain running and service times, control-point residence, and later dispatch availability remain outside this scalar accounting identity. It also does not compare total holding with an unconstrained policy that would visit different future states. Its practical value is that prediction and execution refer to the same limited resource.

\subsection{Timetable Revision Redistributes Adjacent Targets}
Let $\tau_-$ be the preceding published release. Before a permitted revision, the two adjacent target headways are $p_1-\tau_-$ and $p_2-p_1$. After changing the next release to $\tau_1$, they become $\tau_1-\tau_-$ and $p_2-\tau_1$. Therefore,
\begin{equation}
 (\tau_1-\tau_-)+(p_2-\tau_1)=p_2-\tau_-.
 \label{eq:anchor}
\end{equation}
The law redistributes service within a fixed anchored span. It neither inserts a trip nor increases the nominal fleet. The term ``frequency-informed'' refers to temporal components of the observed signal, rather than an increase in the total scheduled service frequency. As actual releases depend on vehicle availability,~\eqref{eq:anchor} describes published targets and does not force realized headways to attain those values.

This limited intervention also explains why many budget and joint choices can be behaviorally identical. The same state may yield no mutable departure, an unsupported supply estimate, an unchanged rounded release, or a release whose effect is absorbed before dispatch. A critic trained on the mode label then learns across both effective and ineffective interventions. The relevant learning question is whether state information identifies the occasions on which a permitted revision changes passenger outcomes, not whether the actor selects joint often.

\subsection{Identical Cost Rankings Can Produce Different Actions}
Let $\widehat C(h)$ be the predicted passenger cost of a current action, and let $\mathcal F(\tau,B)$ be its feasible set under the published plan and remaining budget. The fast decision is
\begin{equation}
 h^*(\tau,B)\in\arg\min_{h\in\mathcal F(\tau,B)}\widehat C(h).
\end{equation}
Even if a plan change leaves $\widehat C(h_1)<\widehat C(h_2)$ intact, it can switch the selected action by removing $h_1$ from the feasible set. A change in the fast command is thus evidence of operational coupling without, by itself, proving that the new plan improves passenger cost. The matched capture in the experiments provides precisely this kind of observation. It is a useful mechanism test because all retained non-plan inputs can be held equal.

\subsection{Consistent Windows for Regularity Comparisons}
For $n$ included arrival headways with sum $S$ and squared sum $Q$, the population-moment convention gives
\begin{equation}
 \mathrm{CV}^2=\frac{nQ}{S^2}-1,\qquad n>0,\ S>0.
 \label{eq:cv_moments}
\end{equation}
The fast model extends an observed prefix with a predicted suffix. A physical short continuation may end before some of the suffix events occur. Comparing its CV$^2$ with a complete model rollout then changes $n$, $S$, and $Q$ together. Even the sign of an action difference can change when a single boundary event is included. Our diagnostic recomputes all three quantities on a common observation window before interpreting a sign disagreement as a prediction failure. Window alignment does not correct inaccurate event times; it makes the comparison address the same set of temporal obligations.

Together, these properties locate the method's guarantees at the interface: cumulative command feasibility, anchored plan semantics, and well-defined comparisons. Realized waiting-time improvement, vehicle readiness, and generalization remain empirical questions.
